\subsection{Squarefree certificates and conditional prime-primary linear rigidity}
\label{sec:linear-collision-rigidity}

The total degree of a critical Mahler certificate can have logarithmic
growth because multiplicities accumulate along power orbits.  The number of
distinct zeros and poles has a sharper bound.  For a reduced normalized
rational function
\[
 R=\frac AB,
 \qquad (A,B)=1,
 \qquad A(0)=B(0)=1,
\]
put
\[
 \rho(R):=\deg\rad(AB).
\]
Thus \(\rho(R)\) is the number of distinct non-zero zeros and poles over
\(\overline{\QQ}\), without multiplicity.

\begin{theorem}[Sharp squarefree Mahler bound]
\label{thm:squarefree-mahler-bound}
Let \(p\ge2\), and suppose that
\begin{equation}\label{eq:squarefree-mahler-coboundary}
 H(z)=\frac{P_0(z)}{P_1(z)}
 =\frac{R(z^p)}{R(z)^p},
\end{equation}
where \(P_0,P_1\in\QQ[z]\) are coprime,
\(P_0(0)=P_1(0)=1\), and \(R\in\QQ(z)^\times\) satisfies \(R(0)=1\).
If
\[
 D=\deg P_0+\deg P_1,
\]
then
\begin{equation}\label{eq:sharp-squarefree-mahler-bound}
 2(p-1)\rho(R)\le D.
\end{equation}
The constant \(2(p-1)\) is sharp for every \(p\ge2\).
\end{theorem}

\begin{proof}
For \(\alpha\in\overline{\QQ}^{\times}\), write
\[
 e(\alpha)=\operatorname{ord}_{z=\alpha}H,
 \qquad
 r(\alpha)=\operatorname{ord}_{z=\alpha}R.
\]
Let \(S=\operatorname{supp}r\) and \(s=\abs S=\rho(R)\).  Since the input
in \eqref{eq:squarefree-mahler-coboundary} is reduced and has constant term
one,
\begin{equation}\label{eq:squarefree-input-l1-degree}
 D=\sum_{\alpha\in\overline{\QQ}^{\times}}\abs{e(\alpha)}.
\end{equation}
Taking divisors gives the identity already used in
\eqref{eq:p-mahler-divisor-recurrence}:
\begin{equation}\label{eq:squarefree-divisor-recurrence}
 e(\alpha)=r(\alpha^p)-p r(\alpha).
\end{equation}

For \(\beta\in S\), put
\[
 d_\beta=\#\{\alpha\in S:\alpha^p=\beta\}.
\]
Define a finitely supported function
\(\varphi:\overline{\QQ}^{\times}\to\{-1,0,1\}\) by
\[
 \varphi(\alpha)=
 \begin{cases}
  -\operatorname{sgn}r(\alpha),&\alpha\in S,\\
  \operatorname{sgn}r(\alpha^p),&\alpha\notin S,
       \ \alpha^p\in S,\\
  0,&\text{otherwise}.
 \end{cases}
\]
Equations \eqref{eq:squarefree-input-l1-degree} and
\eqref{eq:squarefree-divisor-recurrence} imply
\begin{align}
 D
 &\ge \sum_\alpha e(\alpha)\varphi(\alpha)\notag\\
 &=\sum_{\beta\in S}r(\beta)
   \left(\sum_{\alpha^p=\beta}\varphi(\alpha)-p\varphi(\beta)\right).
 \label{eq:squarefree-dual-pairing}
\end{align}
Fix \(\beta\in S\).  Exactly \(p-d_\beta\) of its \(p\) distinct
\(p\)-th roots lie outside \(S\), and each contributes
\(\abs{r(\beta)}\) to the grouped summand in
\eqref{eq:squarefree-dual-pairing}.  Each of the \(d_\beta\) internal roots
contributes at least \(-\abs{r(\beta)}\), while
\[
 -p r(\beta)\varphi(\beta)=p\abs{r(\beta)}.
\]
The total contribution belonging to \(\beta\) is therefore at least
\(2(p-d_\beta)\abs{r(\beta)}\), and hence at least
\(2(p-d_\beta)\).  Consequently
\[
 D\ge2\sum_{\beta\in S}(p-d_\beta)
 =2\left(ps-\sum_{\beta\in S}d_\beta\right).
\]
Now
\[
 \sum_{\beta\in S}d_\beta
 =\#\{\alpha\in S:\alpha^p\in S\}\le s,
\]
which proves \eqref{eq:sharp-squarefree-mahler-bound}.

For sharpness, take \(R(z)=1-z^q\), where \(q\ge1\).  Then
\(\rho(R)=q\), and
\[
 \frac{R(z^p)}{R(z)^p}
 =\frac{1+z^q+\cdots+z^{(p-1)q}}{(1-z^q)^{p-1}}
\]
is reduced.  Its numerator and denominator both have degree
\((p-1)q\), so equality holds in
\eqref{eq:sharp-squarefree-mahler-bound}.
\end{proof}

The squarefree estimate controls collisions and a Taylor jet at the same
time.

\begin{theorem}[Collision--jet inequality]
\label{thm:collision-jet-inequality}
Under the hypotheses of \Cref{thm:squarefree-mahler-bound}, suppose in
addition that, for some \(x>0\),
\[
 R(t)>0\qquad(0\le t\le x),
\]
and that \(R\ne1\).  Put
\[
 \nu=\operatorname{ord}_{z=0}(R(z)-1),
 \qquad
 c_x(R)=\#\{y\in(0,x):R(y)=1\}.
\]
Then
\begin{equation}\label{eq:collision-jet-bound}
 c_x(R)+\nu\le\rho(R)
 \le\left\lfloor\frac{D}{2(p-1)}\right\rfloor.
\end{equation}
Moreover
\[
 \operatorname{ord}_{z=0}(H(z)-1)=\nu.
\]
\end{theorem}

\begin{proof}
The real function \(f(t)=\log R(t)\) has a zero of order \(\nu\) at zero.
If
\[
 0<y_1<\cdots<y_k<x
\]
are all its other zeros, then \(k=c_x(R)\).  The derivative
\(f'=R'/R\) has a zero of order at least \(\nu-1\) at zero, and Rolle's
theorem supplies a further zero in each of
\[
 (0,y_1),(y_1,y_2),\ldots,(y_{k-1},y_k).
\]
On the other hand, logarithmic differentiation gives
\[
 \frac{R'(z)}{R(z)}
 =\sum_{\alpha\in S}\frac{r(\alpha)}{z-\alpha}
 =\frac{N(z)}{\prod_{\alpha\in S}(z-\alpha)},
\]
where \(N\ne0\) and \(\deg N\le\abs S-1=\rho(R)-1\).  Counting the
zeros of \(N\), with multiplicity, yields
\[
 (\nu-1)+k\le\rho(R)-1.
\]
This proves the first inequality in \eqref{eq:collision-jet-bound}; the
second is \Cref{thm:squarefree-mahler-bound}.

Finally, if \(R(z)=1+az^\nu+O(z^{\nu+1})\), with \(a\ne0\), then
\[
 \frac{R(z^p)}{R(z)^p}=1-pa z^\nu+O(z^{\nu+1}),
\]
because \(p\nu>\nu\).  This proves the order identity.
\end{proof}

We now impose only a relative symmetry on a pair of extensions.  Let
\(\ell\) be prime, let \(G\) be a finite abelian \(\ell\)-group, and set
\[
 P_{\chi,\tau}(z)=\det(I-zA_{\chi,\tau}),
 \qquad
 H_\chi(z)=\frac{P_{\chi,\tau}(z)}{P_{\chi,\tau'}(z)}.
\]
The pair \((\tau,\tau')\) is \emph{relatively unit-Adams invariant} if
\begin{equation}\label{eq:relative-unit-adams}
 H_{\chi^u}(z)=H_\chi(z)
 \qquad
 (\chi\in\widehat G,\ (u,\ell)=1).
\end{equation}
This is weaker than requiring the two determinant families to be separately
unit-Adams invariant.

\begin{lemma}[Prime-primary Adams--M\"obius collapse]
\label{lem:prime-primary-adams-collapse}
Assume \eqref{eq:relative-unit-adams}, and put
\[
 \Delta_\chi(z)=F_\chi^\tau(z)-F_\chi^{\tau'}(z)
 =-\sum_{g\in G}\chi(g)
  \bigl(\mathcal E_g^\tau(z)-\mathcal E_g^{\tau'}(z)\bigr).
\]
Then \(H_\chi\in\QQ(z)\), and on the common open Perron disk
\begin{equation}\label{eq:prime-primary-adams-collapse}
 \Delta_\chi(z)
 =-\sum_{j\ge0}\ell^{-j}\Log_0H_\chi(z^{\ell^j})
  +\sum_{j\ge1}\ell^{-j}\Log_0H_{\chi^\ell}(z^{\ell^j}).
\end{equation}
\end{lemma}

\begin{proof}
The abelian form of \eqref{eq:fixed-determinant-open-disc} gives
\[
 \Delta_\chi(z)
 =-\sum_{k\ge1}\frac1k\sum_{d\mid k}\mu(d)
       \Log_0H_{\chi^d}(z^k).
\]
Write \(k=\ell^j q\), where \((q,\ell)=1\).  The squarefree divisors with
non-zero M\"obius value are \(d\mid q\), and, when \(j\ge1\), also
\(\ell d\).  Relative unit-Adams invariance gives
\[
 H_{\chi^d}=H_\chi,
 \qquad
 H_{\chi^{\ell d}}=H_{\chi^\ell}.
\]
The factor \(\sum_{d\mid q}\mu(d)\) vanishes unless \(q=1\), leaving
exactly \eqref{eq:prime-primary-adams-collapse}.

If \(K=\QQ(\chi(G))\), every element of \(\operatorname{Gal}(K/\QQ)\)
sends \(\chi\) to \(\chi^u\) for some \((u,\ell)=1\), and it sends
\(H_\chi\) to \(H_{\chi^u}\).  Thus \eqref{eq:relative-unit-adams} fixes
\(H_\chi\) under the full Galois group, proving
\(H_\chi\in\QQ(z)\).
\end{proof}

\begin{theorem}[Conditional prime-primary collision--jet rigidity]
\label{thm:prime-primary-collision-jet-rigidity}
Let \(\ell\) be prime and let \(G\) be a finite abelian \(\ell\)-group.
Let \(\tau,\tau'\) be one-step \(G\)-cocycles over primitive non-negative
integral base matrices \(A,A'\), of sizes \(v,v'\), with Perron roots
\(\lambda,\lambda'>1\).  Put
\[
 x=\min\{\lambda^{-1},(\lambda')^{-1}\},
\]
and assume relative unit-Adams invariance
\eqref{eq:relative-unit-adams}.  Suppose that, for integers \(L\ge0\) and
\(K\ge1\),
\begin{align}
 p_{n,g}(\tau)&=p_{n,g}(\tau')
 &&(1\le n\le L,\ g\in G),
 \label{eq:prime-primary-initial-jet}\\
 \mathcal E_g^\tau(y_i)&=\mathcal E_g^{\tau'}(y_i)
 &&(g\in G,\ 1\le i\le K),
 \label{eq:prime-primary-profile-samples}
\end{align}
at distinct radii \(0<y_i<x\), at least one of which is algebraic.  Assume
also the following statement, denoted \textup{(KN85)} and attributed to
Keiji Nishioka's 1985 paper, whose bibliographic record is verified but whose
original printed statement has not been verified:
\begin{quote}
For every integer \(p\ge2\) and every \(H\in\CC(z)^\times\), every
convergent Laurent-series germ \(F\in\CC((z))\) that is algebraic over
\(\CC(z)\) and satisfies
\begin{equation}\label{eq:kn85-explicit-assumption}
 F(z^p)=\frac{F(z)^p}{H(z)}
\end{equation}
belongs to \(\CC(z)\),
\end{quote}
If
\begin{equation}\label{eq:prime-primary-cross-base-threshold}
 K+L\ge\max\{v,v'\},
\end{equation}
then
\[
 p_{n,g}(\tau)=p_{n,g}(\tau')
 \qquad(n\ge1,\ g\in G).
\]
If the base determinants are already equal, it is enough to assume
\begin{equation}\label{eq:prime-primary-fixed-determinant-threshold}
 K+L\ge
 \left\lfloor\frac{v+v'}{2(\ell-1)}\right\rfloor.
\end{equation}
\end{theorem}

\begin{proof}
Summing \eqref{eq:prime-primary-profile-samples} over \(g\) gives
\[
 \log\det(I-y_iA)=\log\det(I-y_iA')
 \qquad(1\le i\le K).
\]
Both determinants are positive on \((0,x)\), so their values agree at the
\(y_i\).  Likewise, summing \eqref{eq:prime-primary-initial-jet} over \(g\)
and using the primitive-orbit product gives
\[
 \det(I-zA)-\det(I-zA')=O(z^{L+1}).
\]
The polynomial on the left has degree at most \(\max\{v,v'\}\).  Under
\eqref{eq:prime-primary-cross-base-threshold}, its zero of multiplicity at
least \(L+1\) at zero and its \(K\) distinct positive zeros force it to
vanish identically.  Thus in either case under consideration the base
determinants agree and
\begin{equation}\label{eq:prime-primary-trivial-ratio}
 H_{\mathbf1}=1.
\end{equation}

Fourier transformation of \eqref{eq:prime-primary-profile-samples} gives
\begin{equation}\label{eq:prime-primary-fourier-samples}
 \Delta_\chi(y_i)=0
 \qquad(\chi\in\widehat G,\ 1\le i\le K).
\end{equation}
We prove \(H_\chi=1\) by induction on the order of \(\chi\).  The trivial
case is \eqref{eq:prime-primary-trivial-ratio}.  Let \(\chi\) have order
\(\ell^a\) and suppose that all lower-order character ratios are one.  Then
\(H_{\chi^\ell}=1\), and
\Cref{lem:prime-primary-adams-collapse} reduces to
\begin{equation}\label{eq:prime-primary-one-orbit}
 \Delta_\chi(z)
 =-\sum_{j\ge0}\ell^{-j}\Log_0H_\chi(z^{\ell^j}).
\end{equation}

By \Cref{lem:prime-primary-adams-collapse}, \(H_\chi\in\QQ(z)\).  The
Perron-weighted domination in \Cref{lem:perron-weighted-domination} makes its
numerator and denominator zero-free on \(\abs z<x\).  Since
\(H_\chi(0)=1\) and \(H_\chi\) is real, it is positive on \((0,x)\).
Choose an algebraic sampled radius \(y_0\).  If \(H_\chi\ne1\), then
\eqref{eq:prime-primary-fourier-samples} and
\eqref{eq:prime-primary-one-orbit} give
\begin{equation}\label{eq:prime-primary-algebraic-product}
 \prod_{j\ge0}H_\chi(y_0^{\ell^j})^{\ell^{-j}}=1.
\end{equation}
Define the convergent germ
\[
 F_\chi(z)=\prod_{j\ge0}H_\chi(z^{\ell^j})^{\ell^{-j}}.
\]
It belongs to \(1+z\QQ[[z]]\) and satisfies
\begin{equation}\label{eq:prime-primary-critical-equation}
 F_\chi(z^\ell)=\frac{F_\chi(z)^\ell}{H_\chi(z)}.
\end{equation}
The proof of \Cref{thm:general-rational-p-mahler-lifting}, through its
application of Kumiko Nishioka's 1982 theorem, applies verbatim: one has
\(0<y_0<x<1\), a radius \(r\in(y_0,x)\) gives the required zero- and
pole-free disk, and that proof establishes the coefficient-height and
common-denominator bounds for every rational input.  Hence
\eqref{eq:prime-primary-algebraic-product}, which says \(F_\chi(y_0)=1\),
forces \(F_\chi\) to be algebraic over \(\CC(z)\).  Assumption
\textup{(KN85)} applied to \eqref{eq:prime-primary-critical-equation} now
gives \(F_\chi\in\CC(z)\).  Since its Taylor coefficients are rational,
coefficientwise Galois invariance gives \(F_\chi\in\QQ(z)\).  Therefore
\(R_\chi=F_\chi^{-1}\) is the unique normalized element of
\(\QQ(z)^\times\) such that
\begin{equation}\label{eq:prime-primary-rational-certificate}
 H_\chi(z)=\frac{R_\chi(z^\ell)}{R_\chi(z)^\ell},
 \qquad R_\chi(y_0)=1.
\end{equation}
The positive product defining the lift also shows that
\(R_\chi(t)>0\) for \(0\le t<x\).

Substitution of \eqref{eq:prime-primary-rational-certificate} into
\eqref{eq:prime-primary-one-orbit} telescopes to
\begin{equation}\label{eq:prime-primary-rational-telescope}
 \Delta_\chi(z)=\ell\log R_\chi(z).
\end{equation}
Thus \eqref{eq:prime-primary-fourier-samples} gives
\(R_\chi(y_i)=1\) for every \(i\).  Moreover,
\eqref{eq:prime-primary-initial-jet} gives
\(\Delta_\chi(z)=O(z^{L+1})\), so
\begin{equation}\label{eq:prime-primary-certificate-jet}
 \nu_\chi:=\operatorname{ord}_{z=0}(R_\chi(z)-1)\ge L+1.
\end{equation}

Write \(H_\chi=P_{0,\chi}/P_{1,\chi}\) in reduced form and put
\(D_\chi=\deg P_{0,\chi}+\deg P_{1,\chi}\).  Cancellation from the two
original determinants can only lower degrees, hence
\begin{equation}\label{eq:prime-primary-input-degree}
 D_\chi\le v+v'.
\end{equation}
Applying \Cref{thm:collision-jet-inequality} to
\eqref{eq:prime-primary-rational-certificate} gives
\begin{equation}\label{eq:prime-primary-final-count}
 K+L+1\le K+\nu_\chi
 \le\rho(R_\chi)
 \le\left\lfloor\frac{D_\chi}{2(\ell-1)}\right\rfloor.
\end{equation}
The last member is at most \(\max\{v,v'\}\), contradicting
\eqref{eq:prime-primary-cross-base-threshold}; when the base determinants were
known in advance, \eqref{eq:prime-primary-input-degree} instead makes
\eqref{eq:prime-primary-final-count} contradict
\eqref{eq:prime-primary-fixed-determinant-threshold}.  Therefore
\(H_\chi=1\), closing the induction.

All character determinants now agree.  The comparison dictionary in
\Cref{prop:comparison-dictionary}, followed by Fourier inversion on \(G\),
recovers every primitive length--element count.
\end{proof}

\begin{corollary}[Conditional linear radial determination]
\label{cor:prime-primary-linear-radial-determination}
Assume \textup{(KN85)}: for every \(p\ge2\) and
\(H\in\CC(z)^\times\), every convergent Laurent-series germ
\(F\in\CC((z))\) algebraic over \(\CC(z)\) and satisfying
\(F(z^p)=F(z)^p/H(z)\) belongs to \(\CC(z)\).  Under the remaining
hypotheses of \Cref{thm:prime-primary-collision-jet-rigidity}, with \(L=0\),
\(\max\{v,v'\}\) full profile vectors, one algebraically anchored, determine
all primitive length--element data.  In particular, \(V\) radial locations
suffice for bases of size at most \(V\).
\end{corollary}

\begin{corollary}[Conditional binary linear sampling bound]
\label{cor:exact-binary-radial-sampling-order}
Assume \textup{(KN85)}: for every \(p\ge2\) and
\(H\in\CC(z)^\times\), every convergent Laurent-series germ
\(F\in\CC((z))\) algebraic over \(\CC(z)\) and satisfying
\(F(z^p)=F(z)^p/H(z)\) belongs to \(\CC(z)\).  Then, for \(V\ge6\),
the universal binary radial determination number from
\Cref{cor:linear-radial-sampling-lower-bound} satisfies
\begin{equation}\label{eq:exact-binary-radial-sampling-order}
 \left\lfloor\frac{V-2}{4}\right\rfloor+1
 \le\mathsf N_{C_2}(V)\le V.
\end{equation}
Thus \textup{(KN85)} would imply
\(\mathsf N_{C_2}(V)=\Theta(V)\).  Unconditionally, the displayed lower
bound remains valid.
\end{corollary}

\subsection{Odd-prime multi-collisions with delayed discrepancy}
\label{sec:odd-prime-delayed-multicollisions}

The following unconditional construction gives odd-prime examples with
linearly many collisions and a delayed first discrepancy.

\begin{theorem}[Odd-prime multi-collisions with delayed first discrepancy]
\label{thm:odd-prime-delayed-multicollisions}
Let \(\ell\ge3\) be prime and \(m\ge1\).  There are two one-step
\(C_\ell\)-cocycles \(\tau_m,\tau_m'\) over the same positive primitive base
with
\[
 V_m=\ell(2m+1)
\]
vertices such that:
\begin{enumerate}[label=\textup{(\roman*)},leftmargin=*]
 \item every non-trivial character determinant is invariant under all unit
 Adams operations;
 \item both cocycles have strict twisted spectral gap;
 \item their complete element-profile vectors agree at \(m\) distinct
 rational radii in the open Perron interval;
 \item their primitive length--element data agree for every \(n\le m\);
 \item their first disagreement is at length \(m+1\), with the explicit
 vector in \eqref{eq:odd-prime-first-discrepancy} below.
\end{enumerate}
\end{theorem}

\begin{proof}
Put
\begin{equation}\label{eq:odd-prime-construction-parameters}
 a=\ell^\ell,
 \qquad K=4\ell^2,
 \qquad s_i=a\bigl(K(m+1)+i\bigr)\quad(1\le i\le m),
\end{equation}
and let \(q=2m+1\).  Index a \(q\)-square matrix \(C_m\) by
\[
 u_0,\ldots,u_m,w_1,\ldots,w_m.
\]
Its only non-zero entries are
\begin{align*}
 (C_m)_{u_{i-1},u_i}&=-a,&
 (C_m)_{u_{i-1},w_i}&=s_i,&
 (C_m)_{w_i,u_i}&=s_i &&(1\le i\le m),\\
 (C_m)_{u_m,u_0}&=-a.
\end{align*}
As in the staged-cycle calculation in
\Cref{thm:realizable-linear-multicollisions}, every non-trivial cycle uses
the closing edge and chooses at each stage either the direct edge or the
two-edge detour.  Hence
\begin{equation}\label{eq:odd-prime-staged-determinant}
 Q_m(z):=\det(I-zC_m)
 =1+az\prod_{i=1}^m(-az+s_i^2z^2).
\end{equation}
All entries of \(C_m\) are divisible by \(\ell^\ell\).

Define two \(\ell q\)-square matrices.  The first is cyclic in
\(q\)-square blocks, with
\begin{equation}\label{eq:odd-prime-cyclic-blocks}
 (B_m)_{0,1}=\frac{C_m}{\ell^{\ell-1}},
 \qquad
 (B_m)_{j,j+1}=\ell I_q\quad(1\le j\le\ell-2),
 \qquad
 (B_m)_{\ell-1,0}=\ell I_q,
\end{equation}
and all other blocks zero.  Put
\begin{equation}\label{eq:odd-prime-direct-sum-block}
 B_m'=C_m^{\oplus\ell}.
\end{equation}
The blocks in \eqref{eq:odd-prime-cyclic-blocks} are integral, and their
product around the block cycle is \(C_m\).  The block-companion determinant
identity and the direct-sum identity give
\begin{equation}\label{eq:odd-prime-block-determinants}
 \det(I-zB_m)=Q_m(z^\ell),
 \qquad
 \det(I-zB_m')=Q_m(z)^\ell.
\end{equation}
Every entry of both matrices is divisible by \(\ell\).

Let
\begin{equation}\label{eq:odd-prime-common-base}
 S=\ell s_m,
 \qquad
 A_m=S J_{V_m},
\end{equation}
where \(J_{V_m}\) is the all-ones matrix.  For
\(M\in\{B_m,B_m'\}\), define the edge-count matrices
\begin{equation}\label{eq:odd-prime-label-counts}
 N_0(M)=\frac{A_m+(\ell-1)M}{\ell},
 \qquad
 N_g(M)=\frac{A_m-M}{\ell}\quad(g\in C_\ell\setminus\{0\}).
\end{equation}
They are integral.  The choice of \(S\) gives
\[
 S>\max_{i,j}\abs{M_{ij}},
 \qquad
 S>(\ell-1)\max_{i,j}(-M_{ij}),
\]
so they are non-negative.  Moreover
\[
 N_0(M)+\sum_{g\ne0}N_g(M)=A_m.
\]
For any non-trivial character \(\chi\) of \(C_\ell\),
\(\sum_{g\ne0}\chi(g)=-1\), and therefore its twisted block is
\[
 N_0(M)+\sum_{g\ne0}\chi(g)N_g(M)=N_0(M)-N_g(M)=M.
\]
Thus every non-trivial character block is \(B_m\) for \(\tau_m\) and
\(B_m'\) for \(\tau_m'\), proving unit-Adams invariance.  Since \(A_m\) is
positive and strictly dominates the absolute values of both twisted blocks,
strict Perron--Frobenius monotonicity gives
\[
 \rad(B_m),\rad(B_m')<\rad(A_m).
\]
The common Perron root is
\begin{equation}\label{eq:odd-prime-perron-root}
 \lambda_m=SV_m=\ell^2(2m+1)s_m.
\end{equation}

Set
\begin{equation}\label{eq:odd-prime-collision-radii}
 y_{m,i}=\frac{a}{s_i^2}\qquad(1\le i\le m).
\end{equation}
These numbers are positive, rational, and distinct.  They lie in the open
Perron interval.  Indeed,
\[
 s_1^2\ge a^2K^2(m+1)^2,
 \qquad
 a\lambda_m\le2a^2\ell^2(K+1)(m+1)^2,
\]
and
\[
 K^2=16\ell^4>2\ell^2(K+1)=8\ell^4+2\ell^2.
\]
Thus \(s_1^2>a\lambda_m\), and every
\(y_{m,i}<\lambda_m^{-1}\).  The \(i\)-th factor in
\eqref{eq:odd-prime-staged-determinant} vanishes at \(y_{m,i}\), so
\begin{equation}\label{eq:odd-prime-q-collisions}
 Q_m(y_{m,i})=1.
\end{equation}

For every non-trivial character, the determinant ratio is
\begin{equation}\label{eq:odd-prime-mahler-ratio}
 H_m(z)=\frac{Q_m(z^\ell)}{Q_m(z)^\ell}.
\end{equation}
The collapse in \Cref{lem:prime-primary-adams-collapse}, or direct
telescoping, gives
\[
 F_\chi^{\tau_m}(z)-F_\chi^{\tau_m'}(z)
 =\ell\log Q_m(z).
\]
Equation \eqref{eq:odd-prime-q-collisions} makes every non-trivial Fourier
profile difference vanish at all \(m\) radii.  The trivial channel vanishes
because the base is common, so Fourier inversion proves equality of the
complete element-profile vectors there.

Finally,
\begin{equation}\label{eq:odd-prime-q-leading-term}
 Q_m(z)=1+c_mz^{m+1}+O(z^{m+2}),
 \qquad c_m=(-1)^m a^{m+1}.
\end{equation}
Since \(\ell(m+1)>m+1\), equations
\eqref{eq:odd-prime-mahler-ratio} and
\eqref{eq:odd-prime-q-leading-term} give
\[
 \log H_m(z)=-\ell c_mz^{m+1}+O(z^{m+2}).
\]
Thus every twisted trace agrees through length \(m\), while
\[
 \Tr(B_m^{m+1})-\Tr((B_m')^{m+1})
 =\ell(m+1)c_m.
\]
The trivial trace difference is zero.  Fourier inversion followed by the
triangular primitive-orbit inversion proves equality of every primitive
length--element count for \(n\le m\).  At the first differing length the
lower-length terms have already cancelled, and Fourier inversion gives
\begin{equation}\label{eq:odd-prime-first-discrepancy}
 \begin{aligned}
 p_{m+1,0}(\tau_m)-p_{m+1,0}(\tau_m')&=(\ell-1)c_m,\\
 p_{m+1,g}(\tau_m)-p_{m+1,g}(\tau_m')&=-c_m
 \qquad(g\ne0).
 \end{aligned}
\end{equation}
This also proves that the two periodic-data sets are different.
\end{proof}
